\section{Missing Proofs for New Implications}
\label{sec:new-impls-extra}

\thmRecall{rthm:impl:mxs-to-ef1-n2}
\begin{proof}
Suppose $A$ is MXS-fair to agent 1 but not EF1-fair to her.
Then agent 1 envies agent 2 in $A$, so $v_1(A_1) < v_1(A_2)$.
Let $B$ be agent 1's MXS-certificate for $A$. Then $v_1(B_1) \le v_1(A_1)$.
Moreover, $v_1(A_2) = v_1([m]) - v_1(A_1) \le v_1([m]) - v_1(B_1) = v_1(B_2)$.
Hence, we get $v_1(B_1) \le v_1(A_1) < v_1(A_2) \le v_1(B_2)$.

Let $G \defeq \{g \in [m]: v_1(g) > 0\}$ and $C \defeq \{c \in [m]: v_1(c) < 0\}$.
Let $\max(\emptyset) \defeq -\infty$ and $\min(\emptyset) \defeq \infty$.

Since agent 1 is EFX-satisfied by $B$ and not EF1-satisfied by $A$,
for every $\ghat \in A_2$, we get
\begin{align*}
& \frac{v_1(A_2) - v_1(\ghat)}{w_2} > \frac{v_1(A_1)}{w_1}
\ge \frac{v_1(B_1)}{w_1}
\\ &\ge \frac{1}{w_2}\left(v_1(B_2) - \min_{g \in B_2 \cap G} v_1(g)\right)
\\ &\ge \frac{1}{w_2}\left(v_1(A_2) - \min_{g \in B_2 \cap G} v_1(g)\right).
\end{align*}
Hence, for every $\ghat \in A_2$, we get $v_1(\ghat) < \min_{g \in B_2 \cap G} v_1(g)$.
Hence, $A_2 \cap G$ and $B_2 \cap G$ are disjoint, so $A_2 \cap G \subseteq B_1 \cap G$.
Let $d_i \defeq -v_i$ for all $i$. Then for every $\chat \in A_1$, we get
\begin{align*}
& \frac{d_1(A_1) - d_1(\chat)}{w_1} > \frac{d_1(A_2)}{w_2}
\ge \frac{d_2(B_2)}{w_2}
\\ &\ge \frac{1}{w_1}\left(d_1(B_1) - \min_{c \in B_1 \cap C} d_1(c)\right)
\\ &\ge \frac{1}{w_1}\left(d_1(A_1) - \min_{c \in B_1 \cap C} d_1(c)\right).
\end{align*}
Hence, for every $\chat \in A_1$, we have $d_1(\chat) < \min_{c \in B_1 \cap C} d_1(c)$.
Hence, $A_1 \cap C$ and $B_1 \cap C$ are disjoint, so $B_1 \cap C \subseteq A_2 \cap C$.
Hence,
\begin{align*}
v_1(A_2) &= v_1(A_2 \cap G) - d_1(A_2 \cap C)
\\ &\le v_1(B_1 \cap G) - d_1(B_1 \cap C) = v_1(B_1),
\end{align*}
which is a contradiction.
Hence, it's not possible for $A$ to be MXS-fair to agent 1 but not EF1-fair to her.
\end{proof}

\thmRecall{rthm:impl:eef1-to-prop1}
\begin{proof}
Suppose agent $i$ is epistemic-EF1-satisfied but not PROP1-satisfied by allocation $X$.
Let $A$ be agent $i$'s epistemic-EF1-certificate for $X$.
Then $A$ is EF1-fair to agent $i$.
If $A$ is PROP1-fair to $i$, then $X$ would also be PROP1-fair to agent $i$,
which is a contradiction. Hence, $A$ is not PROP1-fair to $i$.
Therefore, all of these hold:
\begin{tightenum}
\item $v_i(A_i) < w_iv_i([m])$.
\item $v_i(A_i \cup \{g\}) \le w_iv_i([m])$ for all $g \in [m] \setminus A_i$.
\item $v_i(A_i \setminus \{c\}) \le w_iv_i([m])$ for all $c \in A_i$.
\end{tightenum}

Since $v_i$ is subadditive, there exists $j \in [n] \setminus \{i\}$ such that $v_i(A_j) > w_jv_i([m])$
(otherwise $v_i([m]) \le \sum_{j=1}^n v_i(A_j) < \sum_{j=1}^n w_jv_j([m]) = v_i([m])$). Hence,
\[ \frac{v_i(A_i)}{w_i} < v_i([m]) < \frac{v_i(A_j)}{w_j}. \]
Hence, $i$ envies $j$. But $i$ is EF1-satisfied. Hence,
\\ $\exists c \in A_i$ such that
    $\displaystyle \frac{v_i(A_i \setminus \{c\})}{w_i} \ge \frac{v_i(A_j)}{w_j}$,
\\ or $\exists g \in A_j$ such that
    $\displaystyle \frac{v_i(A_i)}{w_i} \ge \frac{v_i(A_j \setminus \{g\})}{w_j}$.

\textbf{Case 1}: $\exists c \in A_i$ such that $v_i(A_i \setminus \{c\})/w_i \ge v_i(A_j)/w_j$.

Since $i$ is PROP1-unsatisfied and $v_i(A_j) > w_jv_i([m])$, we get
\[ \frac{v_i(A_j)}{w_j} \le \frac{v_i(A_i \setminus \{c\})}{w_i} \le v_i([m]) < \frac{v_i(A_j)}{w_j}, \]
which is a contradiction.
Hence, it's impossible for agent $i$ to be EF1-satisfied but not PROP1-satisfied by $A$ for this case.

\textbf{Case 2}: $\exists g \in A_j$ such that $v_i(A_i)/w_i \ge v_i(A_j \setminus \{g\})/w_j$.

First, we show that this case doesn't occur if all items in $A_j$ are chores.
Since $i$ envies $j$ but not EF1-envies $j$, we get
\begin{equation}
\label{eq:impl:eef1-to-prop1:gpos}
\frac{v_i(g \mid A_j \setminus \{g\})}{w_j}
= \frac{v_i(A_j) - v_i(A_j \setminus \{g\})}{w_j}
\ge \frac{v_i(A_j)}{w_j} - \frac{v_i(A_i)}{w_i} > 0.
\end{equation}
Hence, this case doesn't occur if all items in $A_j$ are chores.

Since $i$ is not PROP1-satisfied, we get
\begin{align*}
& \frac{v_i(A_i) + v_i(g \mid A_i)}{w_i} \le w_i([m]) < \frac{v_i(A_j)}{w_j}
\\ &\implies \frac{v_i(g \mid A_i)}{w_i} < \frac{v_i(A_j)}{w_j} - \frac{v_i(A_i)}{w_i}
    \le \frac{v_i(g \mid A_j \setminus \{g\})}{w_j}.
\end{align*}

Let $v_i$ be additive. Then we get
\[ \frac{v_i(g)}{w_i} < \frac{v_i(g)}{w_j} \implies w_j < w_i. \]
If $w_i \le w_j$ for all $j \in [n] \setminus \{i\}$, we get a contradiction.

Now suppose $v_i$ is additive and $n=2$.
Then the agents are $i$ and $j$. Note that $w_i + w_j = 1$.
Let $\ghat \in \argmax_{t \in A_j} v_i(t)$.
Since $i$ is not PROP1-satisfied, we get
\begin{align*}
& v_i(A_i \cup \{\ghat\}) \le w_iv_i([m]) = w_iv_i(A_i \cup \{\ghat\}) + w_iv_i(A_j \setminus \{\ghat\})
\\ &\implies \frac{v_i(A_i \cup \{\ghat\})}{w_i} \le \frac{v_i(A_j \setminus \{\ghat\})}{w_j}.
\end{align*}
Based on the assumption of Case 2, we get
\[ \frac{v_i(A_i)}{w_i} \ge \frac{v_i(A_j \setminus \{g\})}{w_j} \ge \frac{v_i(A_j \setminus \{\ghat\})}{w_j}. \]
Hence,
\[ v_i(A_i \cup \{\ghat\}) \le \frac{w_i}{w_j}v_i(A_j \setminus \{\ghat\}) \le v_i(A_i). \]
Therefore, $v_i(\ghat) \le 0$.
Since $\ghat$ is the most-valuable item in $A_j$ to $i$, we get that $v_i(g) \le 0$,
which contradicts \eqref{eq:impl:eef1-to-prop1:gpos}.
Hence, it's impossible for agent $i$ to be EF1-satisfied but not PROP1-satisfied by $A$ for this case.
\end{proof}

\thmRecall{rthm:impl:mms-to-mxs}
\begin{proof}
We will show that agent $i$'s leximin $n$-partition is her MXS-certificate for $A$,
which would prove that $A$ is MXS-fair to agent $i$.

Without loss of generality, assume $i = 1$.
Let $P$ be a leximin $n$-partition of $v_1$ (c.f.~\cref{defn:leximin})
such that $v_1(P_1) \le v_1(P_2) \le \ldots \le v_1(P_n)$.
Then $v_1(P_1) = \MMS_{v_1}^n([m]) \le v_1(A_1)$.

For any agent $j \ge 2$, the allocation $(P_1, P_j)$ is leximin for
the instance $\Icalhat \defeq (\{i, j\}, P_1 \cup P_j, (v_1, v_j), (1/2, 1/2))$.
Hence, agent 1 is MMS-satisfied by $(P_1, P_j)$ in $\Icalhat$.
Since either all items are goods to agent 1 or $v_1$ is additive,
by \cref{thm:impl:mms-to-efx-n2}, agent 1 is EFX-satisfied by $(P_1, P_j)$ in $\Icalhat$.

On considering all values of $j$, we get that agent 1 is EFX-satisfied by $P$ in $\Ical$.
Hence, $P$ is agent 1's MXS-certificate for $A$.
\end{proof}

\thmRecall{rthm:impl:eefx-to-propx}
\begin{proof}
Suppose $A$ is epistemic-EFX-fair to agent $i$ but not PROPx-fair to $i$.
Let $B$ be agent $i$'s epistemic-EFX-certificate for $A$.
Then $B$ is also PROPx-unfair to $i$.

Since $i$ is EFX-satisfied with $B$, we get that for all $j \in [n] \setminus \{i\}$,
\[ \min_{S \in \Scal} \frac{v_i(B_i \setminus S)}{w_i} \ge \frac{v_i(B_j)}{w_j}, \]
where $\Scal \defeq \{S \subseteq B_i: v_i(S \mid B_i \setminus S) < 0\}$.
Add these inequalities for all $j$, weighting each by $w_j$, to get
\[ \min_{S \in \Scal} \frac{v_i(B_i \setminus \{c\})}{w_i} > \sum_{j=1}^n v_i(B_j) \ge v_i([m]), \]
which implies that $B$ is PROPx-fair, a contradiction.
Hence, there can't be an allocation $A$ that is epistemic-EFX-fair to $i$
but not PROPx-fair to $i$.
\end{proof}

\thmRecall{rthm:impl:efx-to-propm}
\begin{proof}
Suppose agent $i$ is EFX-satisfied but not PROPm-satisfied by allocation $A$.
Then for some agent $\jhat \in [n] \setminus \{i\}$, we have
\[ \frac{v_i(A_i)}{w_i} < v_i([m]) < \frac{v_i(A_{\jhat})}{w_{\jhat}}. \]

Since $i$ doesn't EFX-envy $\jhat$, for all $S \subseteq A_i$ such that
$v_i(S \mid A_i \setminus S) < 0$, we get
\[ \frac{v_i(A_i \setminus S)}{w_i} \ge \frac{v_i(A_{\jhat})}{w_{\jhat}} > v_i([m]). \]
Hence, condition \ref{item:propm:chores} (chores condition) of PROPm is satisfied.
Since condition \ref{item:propm:goods} (goods condition) of PROPm is not satisfied, we get that
$T \neq \emptyset$ and $v_i(A_i) + \max(T) \le w_iv_i([m])$.
%
For all $j \in T$, we get
\begin{align}
\frac{v_i(A_i)}{w_i} &\ge \frac{\max(\{v_i(A_j \setminus S): S \subseteq A_j \textrm{ and } v_i(S \mid A_i) > 0\})}{w_j}
    \tag{$i$ doesn't EFX-envy $j$}
\\ &= \frac{v_i(A_j)}{w_j} - \frac{\min(\{v_i(S): S \subseteq A_j \textrm{ and } v_i(S) > 0\})}{w_j}
    \notag
\\ &= \frac{v_i(A_j)}{w_j} - \frac{\tau_j}{w_j} \ge \frac{v_i(A_j)}{w_j} - \frac{\max(T)}{w_j}.
    \label{eqn:impl:efx-to-propm:1}
\end{align}

\textbf{Special case 1}: $w_i \le w_j$ for all $j \in [n] \setminus \{i\}$, and $v_i(A_i) \ge 0$.

Let $J \defeq \{j \in [n] \setminus \{i\}: \tau_j > 0\}$.
For all $j \in J$, using \eqref{eqn:impl:efx-to-propm:1}, we get
\[ \frac{v_i(A_i) + \max(T)}{w_i}
\ge \frac{v_i(A_j)}{w_j} - \frac{\max(T)}{w_j} + \frac{\max(T)}{w_i}
\ge \frac{v_i(A_j)}{w_j}. \]

For all $j \in [n] \setminus (J \cup \{i\})$, we have $\tau_j = 0$, so $v_i(A_j) \le 0$. Hence,
\[ \frac{v_i(A_i)}{w_i} \ge 0 \ge \frac{v_i(A_j)}{w_j}. \]
Therefore,
\[ v([m]) = \sum_{j=1}^n w_j\left(\frac{v_i(A_j)}{w_j}\right)
    < \sum_{j=1}^n w_j\left(\frac{v_i(A_i) + \max(T)}{w_i}\right)
    = \frac{v_i(A_i) + \max(T)}{w_i}. \]
Hence, condition 1 of PROPm is satisfied, which is a contradiction.
Hence, it's impossible for $i$ to be EFX-satisfied by $A$ but not PROPm-satisfied by $A$.

\textbf{Special case 2}: $n=2$

Let the two agents be $i$ and $j$. Since $T \neq \emptyset$, we have $\max(T) = \tau_j > 0$.
Since $i$ is not PROPm-satisfied by $A$, we get
\begin{align}
& v_i(A_i) + \tau_j \le w_iv_i([m]) = w_iv_i(A_i) + w_iv_i(A_j)
\\ &\implies w_jv_i(A_i) + \tau_j \le w_iv_i(A_j)
\\ &\implies \frac{v_i(A_i)}{w_i} \le \frac{v_i(A_j)}{w_j} - \frac{\tau_j}{w_iw_j}
    \le \frac{v_i(A_j)}{w_j} - \frac{\tau_j}{w_j}.
\end{align}
Combining this with equation \eqref{eqn:impl:efx-to-propm:1} gives us $\tau_j = 0$, which is a contradiction.
Hence, it's impossible for $i$ to be EFX-satisfied by $A$ but not PROPm-satisfied by $A$.
\end{proof}

\thmRecall{rthm:wmms-vs-knapsack}
\begin{proof}
For any $j \in [n]$, define
\[ S_j \defeq \argmax_{\substack{S \subseteq [m]:\\v_i(S) \le w_jv_i([m])}} \frac{v_i(S)}{w_j}. \]
Then $\beta_i \defeq \max_{j=1}^n v_i(S_j)/w_j$.

Let $\Pi_n([m])$ be the set of all $n$-partitions of $[m]$. Let
\begin{align*}
P &\defeq \argmax_{P \in \Pi_n([m])} \min_{j=1}^n \frac{v_i(P_j)}{w_j},
& k &\defeq \argmin_{j=1}^n \frac{v_i(P_j)}{w_j}.
\end{align*}
Then $\WMMS_i/w_i \defeq v_i(P_k)/w_k$.

By \cref{thm:impl:prop-to-wmms}, we get $v_i(P_k)/w_k \le v_i([m])$. Hence,
\[ \frac{\WMMS_i}{w_i} = \frac{v_i(P_k)}{w_k} \le \frac{v_i(S_k)}{w_k}
    \le \max_{j=1}^n \frac{v_i(S_j)}{w_j} = \beta_i. \]

Now let $v_i$ be additive and $n = 2$.
For any $j \in [2]$, let $Q^{(j)}$ be an allocation where $Q^{(j)}_j \defeq S_j$
and $Q^{(j)}_{3-j} \defeq [m] \setminus S_j$.
Then $v_i(Q^{(j)}_j) = v_i(S_j) \le w_jv_i([m])$ and
\[ v_i(Q^{(j)}_{3-j}) = v_i([m]) - v_i(S_j) \ge v_i([m]) - w_jv_i([m]) = w_{3-j}v_i([m]). \]
Hence,
\[ \frac{v_i(Q^{(j)}_j)}{w_j} \le v_i([m]) \le \frac{v_i(Q^{(j)}_{3-j})}{w_{3-j}}. \]
Therefore,
\begin{align*}
\frac{\WMMS_i}{w_i} &= \min\left(\frac{v_i(P_1)}{w_1}, \frac{v_i(P_2)}{w_2}\right)
\\ &\ge \max_{j=1}^n \min\left(\frac{v_i(Q^{(j)}_j)}{w_j}, \frac{v_i(Q^{(j)}_{3-j})}{w_{3-j}}\right)
    \tag{by definition of $P$}
\\ &= \max_{j=1}^n \frac{v_i(S_j)}{w_j} = \beta_i.
\end{align*}
Hence, $\WMMS_i/w_i \ge \beta_i$.
\end{proof}

\thmRecall{rthm:impl:mms-to-aps-n2}
\begin{proof}
By \cref{thm:wmms-vs-knapsack}, we get
\[ \frac{\WMMS_i}{w_i} = \beta_i \defeq \max_{j=1}^n \frac{v_i(S_j)}{w_j},
\quad\textrm{where}\quad
S_j \defeq \argmax_{\substack{S \subseteq [m]:\\ v_i(S) \le w_jv_i([m])}} v_i(S). \]
On setting $p = v_i$ in the primal definition of APS (\cref{defn:aps}), we get
\[ \APS_i \le \max_{\substack{S \subseteq [m]:\\ p(S) \le w_ip([m])}} v_i(S)
    = v_i(S_i) \le w_i\beta_i = \WMMS_i. \]
When entitlements are equal, $\APS_i \ge \pessShare_i \ge \MMS_i$
by \cref{thm:impl:aps-to-pess}.
\end{proof}

\thmRecall{rthm:tribool-rr}
\begin{proof}
Partition $[m]$ into
goods $M_+ \defeq \{g \in [m]: v_i(g) > 0\}$,
chores $M_- \defeq \{c \in [m]: v_i(c) < 0\}$,
and neutral items $M_0 \defeq \{t \in [m]: v_i(t) = 0\}$.
Fuse items $M_0$, $\min(|M_+|, |M_-|)$ goods, and $\min(|M_+|, |M_-|)$ chores
into a single item $h$.
Then we are left with only goods and a neutral item,
or only chores and a neutral item.
Using round-robin, one can allocate items such that
any two bundles differ by at most one item.
\end{proof}

\thmRecall{rthm:impl:tribool:prop}
\begin{proof}
Since $A$ is PROP-fair to $i$, and $v_i(S) \in \mathbb{Z}$ for all $S \subseteq [m]$,
we get $v_i(A_i) \ge \ceil{v_i([m])/n}$.

Construct an allocation $B$ where $B_i = A_i$, and items $[m] \setminus A_i$
are allocated among agents $[n] \setminus \{i\}$ using \cref{thm:tribool-rr} with $f = v_i$.
We will show that $B$ is agent $i$'s epistemic-EF-certificate for $A$.

In $B$, for each agent $j \in [n] \setminus \{i\}$, we have
\[ v_i(B_j) \le \bigceil{\frac{v_i([m] \setminus A_i)}{n-1}}
    \le \bigceil{\frac{v_i([m]) - v_i([m])/n}{n-1}}
    \le \bigceil{\frac{v_i([m])}{n}} \le v_i(A_i). \]
Hence, $i$ doesn't envy anyone in $B$.
Hence, $B$ is agent $i$'s epistemic-EF-certificate for $A$.
\end{proof}

\thmRecall{rthm:impl:tribool:prop1}
\begin{proof}
Partition $[m]$ into
goods $M_+ \defeq \{g \in [m]: v_i(g) > 0\}$,
chores $M_- \defeq \{c \in [m]: v_i(c) < 0\}$,
and neutral items $M_0 \defeq \{t \in [m]: v_i(t) = 0\}$.

\textbf{Case 1}: $A_i$ has all goods and no chores.
\\ Then $v_i(A_i) \ge \max(0, v_i([m]))$.
If $v_i([m]) \ge 0$, then $v_i(A_i) \ge v_i([m]) \ge w_iv_i([m])$,
else $v_i(A_i) \ge 0 \ge w_iv_i([m])$.
Hence, $v_i(A_i) \ge w_iv_i([m])$ and $A$ is PROPx+PROP1.

\textbf{Case 2}: $A_i$ has a chore or some good is outside $A_i$.
\\ Then adding a good to $A_i$ or removing a chore from $A_i$
makes it's value more than $w_iv_i([m])$ iff $v_i(A_i) \ge \floor{w_iv_i([m])}$.
Hence, $A$ is PROP1 iff $A$ is PROPx iff $v_i(A_i) \ge \floor{w_iv_i([m])}$.
\end{proof}

\thmRecall{rthm:impl:tribool:aps}
\begin{proof}
Partition $[m]$ into goods $M_+ \defeq \{g \in [m]: v_i(g) > 0\}$,
chores $M_- \defeq \{c \in [m]: v_i(c) < 0\}$, and
neutral items $M_0 \defeq \{t \in [m]: v_i(t) = 0\}$.

First, set $p(t) = v_i(t)$ for each $t \in [m]$ to get $\APS_i \le w_im_i$.
Since the APS is the value of some bundle, and bundle values can only be integers,
we get $\APS_i \le \floor{w_im_i}$.

Pick an arbitrary price vector $p \in \Delta_m$.
We will construct a set $S$ such that $p(S) \le w_ip([m])$ and $v_i(S) \ge w_iv_i([m])$.
Fuse items $M_0$, $\min(|M_+|, |M_-|)$ goods, and $\min(|M_+|, |M_-|)$ chores
into a single item $h$.
Let $M'_+$ and $M'_-$ be the remaining goods and chores, respectively.
Then $M'_+ = \emptyset$ or $M'_- = \emptyset$.
Using techniques from \cref{thm:aps-optimal-price},
we can assume \wLoG{} that $p_g \ge 0$ for all $g \in M'_+$,
$p_c \le 0$ for all $c \in M'_-$, and $p_h = 0$.

\textbf{Case 1}: $M'_- = \emptyset$.
\\ Let $m_i \defeq |M'_+| = v_i([m])$.
Let $S$ be the cheapest $\floor{w_im_i}$ items in $M'_+$.
Then using \cref{thm:prefix-sum-bound}, we get
\[ p(S) \le \frac{\floor{w_im_i}}{m_i}p(M_+) \le w_ip([m]). \]
Then $S$ is affordable and $v_i(S) = \floor{w_im_i}$.
Hence, $\APS_i \ge \floor{w_im_i}$.

\textbf{Case 2}: $M'_+ = \emptyset$.
\\ Let $m_i \defeq |M'_-| = -v_i([m])$.
Let $S$ be the cheapest $\ceil{w_im_i}$ items in $M'_- \cup \{h\}$.
Then using \cref{thm:prefix-sum-bound}, we get
\[ -p(S) \ge \frac{\ceil{w_im_i}}{m_i}(-p([m])) \ge w_i(-p([m])). \]
Then $S$ is affordable and $-v_i(S) \le \ceil{w_im_i} = -\floor{w_iv_i([m])}$.
Hence, $\APS_i \ge \floor{w_iv_i([m])}$.
\end{proof}

\thmRecall{rthm:impl:tribool:m1s}
\begin{proof}
Allocate items $[m]$ among agents $[n]$ using \cref{thm:tribool-rr} with $f = v_i$.
Then any two bundles differ by a value of at most one.
Hence, $\MMS_i = \floor{m'/n}$ and $\mathrm{M1S}_i \le \floor{m'/n}$,
where $m' \defeq v_i([m])$.

Let $X$ be an allocation where agent $i$ is EF1-satisfied.
Then any two bundles can differ by a value of at most one.
Hence, the smallest value $v_i(X_i)$ can have is $\floor{m'/n}$.
Hence, $\mathrm{M1S}_i \ge \floor{m'/n}$.
\end{proof}

\thmRecall{rthm:impl:tribool:ef1-gaps}
\begin{proof}
Consider any subset $S$ of agents.
On restricting $A$ to $S$, we get an allocation $B$ that is EF1-fair to $i$.
$B$ is also PROP1-fair to $i$ by \cref{thm:impl:eef1-to-prop1},
and APS-fair to $i$ by \cref{thm:impl:tribool:prop1,thm:impl:tribool:aps}.
Hence, $A$ is groupwise-APS-fair to $i$.
\end{proof}

\thmRecall{rthm:impl:tribool:mms-to-eefx}
\begin{proof}
Construct an allocation $B$ where $B_i = A_i$, and items $[m] \setminus A_i$
are allocated among agents $[n] \setminus \{i\}$ using \cref{thm:tribool-rr} with $f = v_i$.
We will show that $B$ is agent $i$'s epistemic-EFX-certificate for $A$.

Let $k \defeq v_i([m])$.
Suppose $v_i(A_i) \ge \floor{k/n} = \ceil{(k+1)/n} - 1$ (c.f.~\cref{thm:ceil-floor}).
Then for any other agent $j \in [n] \setminus \{i\}$, we get
\[ v_i(B_j) \le \bigceil{\frac{v_i([m] \setminus A_i)}{n-1}}
    \le \bigceil{\frac{k - (k+1)/n + 1}{n-1}} = \bigceil{\frac{k+1}{n}}
    \le v_i(B_i) + 1. \]
If $B_j$ contains no goods and $B_i$ contains no chores,
then $v_i(B_i) \ge 0 \ge v_i(B_j)$, so $i$ doesn't envy $j$ in $B$.
Otherwise, transferring a good from $j$ to $i$ or a chore from $i$ to $j$ in $B$
eliminates $i$'s envy towards $j$.
Hence, $B$ is agent $i$'s epistemic-EFX-certificate for $A$.
\end{proof}
